\section{Discussió}\label{sec:discussio}

\subsection{Prova de control}\label{sec:control}

Com que TD3 complet no aterrava amb la configuració A, vam comprovar que no fos per un error d'implementació amb una prova fora del protocol, amb les llavors 0--4. La configuració \emph{td3\_paper} aplica els hiperparàmetres de la taula 3 de \textcite{fujimoto2018}: $\tau = 0{,}005$, taxa d'aprenentatge de $10^{-3}$ a l'actor i als crítics, cap \emph{weight decay}, lots de 100 i 10.000 passos inicials d'accions uniformes. \emph{Sense weight decay} només treu el \emph{weight decay} dels crítics. Amb td3\_paper, la llavor 3 resol l'entorn (269,7 punts i 98~\% d'episodis resolts en una avaluació final de 100 episodis), però les llavors 0, 2 i 4 continuen planejant (entre el 83~\% i el 100~\% d'episodis truncats). Sense \emph{weight decay}, el biaix puja en les cinc llavors (de −4,5 a +26,1 de mitjana), però el 79,8~\% dels episodis finals es trunquen.

\subsection{El planeig}\label{sec:planeig}

A TD3 complet amb la configuració A, el 65~\% de les últimes avaluacions tenen els cinc episodis truncats als 1.000 passos (Taula~\ref{tbl:ablacio}): el mòdul no s'estavella ni arriba a repòs. Com que les xarxes finals de l'estudi principal no es van desar, vam mesurar el comportament amb les de la prova de control (20 episodis per a cadascuna de les 10 execucions). Se'n van truncar 143 de 200. En els últims 200 passos, el mòdul era per sobre del terra el 88~\% del temps (\emph{y} mitjana de 0,38; 0 és l'alçada de la plataforma). Les dues potes tocaven terra l'1,6~\% del temps i el motor principal estava encès en el 73~\% dels passos.

La recompensa de LunarLander fa que aquest comportament sigui estable \parencite{gymnasium}. Cada pas amb el motor principal encès costa fins a 0,3 punts, i planejar fins al límit acaba en un retorn moderadament negatiu (−42,4 de mitjana en aquests episodis), mentre que un aterratge fallit suma −100. Aterrar dona +100, però abans cal passar per estats propers al terra en què una acció dolenta porta al xoc. Si aquests estats apareixen poc a la memòria de repetició i el crític els valora malament, és d'esperar que l'actor, que segueix el gradient de $Q_1$, es quedi planejant. Les dades no permeten verificar aquest mecanisme directament.

La configuració B fa el planeig molt menys freqüent, però no l'elimina: de les 8 execucions de B que resolen menys del 50~\% dels episodis finals, 7 en tenen més de la meitat truncats. Els quatre canvis de B es van aplicar junts i no podem dir quin evita el planeig. td3\_paper, que en comparteix dos ($\tau = 0{,}005$ i cap \emph{weight decay}), només aterra en una de cinc llavors.

\subsection{Què canvia cada mecanisme}\label{sec:mecanismes}

El mínim és l'únic mecanisme amb un efecte clar sobre el crític: sense ell apareix la sobreestimació en les dues configuracions. Treure'l també fa l'aprenentatge més inestable: les caigudes de més de 200 punts entre avaluacions consecutives després dels 50.000 passos passen de 17 a 73 a la configuració A i de 13 a 34 a la B ($p = 0{,}031$ i $p = 0{,}016$, tests exactes sense corregir d'una anàlisi exploratòria). A la configuració B, a més, la sobreestimació va lligada al fracàs: a sense mínim, la correlació de Spearman entre el biaix i la L10 és de −0,90, i les tres llavors amb més biaix són les tres que fallen.

L'efecte del mínim sobre el retorn, en canvi, no és el mateix a les dues configuracions: a A, treure'l és el canvi que més millora el retorn, i a B no el millora. Tot i així, la diferència entre els dos efectes (−120,4 de L10, $p = 0{,}107$) no és significativa amb deu llavors. El mínim, per si sol, tampoc no provoca el planeig: TD3 complet el conserva a la configuració B i només hi fallen 2 de 10 llavors.

El suavitzat no té cap efecte clar en el retorn, i sense ell el biaix només puja lleugerament (+5,3 a B, $p = 0{,}047$ sense corregir). Treure el retard, que duplica les actualitzacions de l'actor, millora la mitjana a A i a B, sense significació: a A la millora depèn de dues llavors, i a B hi aterren 9 de 10 llavors, per 8 de TD3 complet.

\subsection{Limitacions}\label{sec:limitacions}

L'ablació, amb deu llavors, no pot detectar diferències de la mida observada, perquè el resultat de cada variant depèn sobretot de quantes llavors fallen. Les xarxes de l'esquelet usen LayerNorm, i l'enunciat parla de \emph{batch normalization}. El biaix es mesura sobre els estats de cada política i per a $\mu$, mentre que el crític aprèn el valor de la política suavitzada.
